\documentclass[11pt]{article}
\usepackage{amsmath,amssymb,geometry,hyperref}
\geometry{margin=1in}
\newcommand{\C}{\mathcal C}
\begin{document}
\title{A certified four-dimensional second-order anticoherent subspace at twenty qubits}
\author{Catalogue calculation}
\date{30 October 2026}
\maketitle

\paragraph{Result.}
There is a complex four-dimensional second-order anticoherent subspace
of $\operatorname{Sym}^{20}(\mathbb C^2)$.  Together with the general
bound $4k\le N-2$ at $N=20$, this proves
\[
\boxed{k_{\max}(20,2)=4.}
\]
The basis below is defined by isolated real algebraic numbers, not by
unverified decimal approximations or an asserted radical formula.

\section{The explicit algebraic basis}
Let $|D_{20}^{(r)}\rangle$ be normalized Dicke states, and put
$\epsilon=(1,1,0,0)$.  Define four vectors
\[
|u_a\rangle=\sum_{\substack{0\le r\le20\\r\equiv\epsilon_a\pmod2}}
\sqrt{\binom{20}{r}}\,(a_{ar}+i b_{ar})|D_{20}^{(r)}\rangle,
\qquad a=1,2,3,4.
\]
The $84$ real coordinates are ordered by increasing $a$, then increasing
$r$ of the indicated parity, then real and imaginary part.  The supplied
\texttt{basis.json} specifies $12$ coordinates as exact rational numbers
and the remaining $72$ coordinates as the unique common real zero of
the displayed rational quadratic polynomials in a displayed rational
box.  Each coordinate is named, each interval endpoint is rational, and
each Dicke coefficient is explicitly associated with those coordinates.
Thus the definition is finite and exact.

\section{All whole-subspace conditions}
Write $B=(u_1,u_2,u_3,u_4)$, $j=10$, and $m_r=10-r$.  The real system
consists of the nonzero parity-supported entries of
\begin{align*}
B^\dagger B&=I_4,&
B^\dagger J_zB&=0,&
B^\dagger(3J_z^2-110I)B&=0,\\
B^\dagger J_+B&=0,&
B^\dagger J_+J_zB&=0,&
B^\dagger J_+^2B&=0.
\end{align*}
For the first three Hermitian matrices it suffices to include the real
diagonal and the real and imaginary parts above the diagonal.  Parity
leaves eight real equations for each matrix.  Each of $J_+$ and
$J_+J_z$ couples only opposite parities and contributes sixteen real
equations.  $J_+^2$ preserves parity and contributes sixteen more.
There are therefore
\[
3\cdot8+2\cdot16+16=72
\]
real equations.  The omitted entries are identically zero by parity or
are Hermitian conjugates of included entries.  There is no selection of
only a convenient subset of the nonzero equations.

In the above coordinates their coefficients are integers.  In particular,
the matrix element of $J_+^qJ_z^p$ is
\[
\sum_{r=q}^{20}\binom{20}{r}\frac{r!}{(r-q)!}\,m_r^p\,
\overline{z_{a,r-q}}z_{b,r},\qquad z_{a,r}=a_{ar}+ib_{ar}.
\]
For the quadrupole equation, replace $m_r^p$ by $3m_r^2-110$.

These conditions give $B^\dagger J_a B=0$ and
$B^\dagger J_aJ_b B=\frac{110}{3}\delta_{ab}I_4$ for all three
Cartesian directions.  Indeed, $J_+^2$ supplies the two transverse
quadrupoles, $J_+J_z$ together with $J_+$ supplies both mixed
quadrupoles, and the diagonal quadrupole plus the Casimir fixes the
remaining second moments.  Thus every normalized vector in the span,
including arbitrary complex superpositions, is second-order AC.
Equivalently, all sixteen ordered two-qubit cross reductions satisfy
\[
\operatorname{Tr}_{18}(|u_a\rangle\langle u_b|)
=\delta_{ab}\frac{I_{\mathrm{Sym}^2}}{3}.
\]

\section{Exact existence certificate}
After fixing the twelve listed rational coordinates, write the complete
system as $F:\mathbb R^{72}\to\mathbb R^{72}$.  The certificate gives
a rational center $c$, a rational matrix $A$, and radius
$r=10^{-35}$.  Define $G(x)=x-AF(x)$.  All bounds below are checked by
integer arithmetic with common denominators; rounded decimals play no
role in the inequalities.

The certificate recomputes
\[
e=\|I-AJ_F(c)\|_\infty<4\cdot10^{-45},\qquad
\beta=\|AF(c)\|_\infty<5\cdot10^{-61}.
\]
Let $H$ bound the maximum, over equations, of the sum of absolute
values of all Hessian entries with respect to the unknown coordinates.
Since every equation is quadratic, $H$ is an explicit integer obtained
from the coefficient list.  Throughout $\|x-c\|_\infty\le r$,
\[
\|DG(x)\|_\infty
\le e+\|A\|_\infty Hr=:\gamma<3\cdot10^{-27}.
\]
Moreover, the exact computation gives
\[
\beta+\gamma r<5.3\cdot10^{-61}<r.
\]
Consequently $G$ maps the closed rational cube strictly into itself and
is a strict contraction.  Banach's theorem gives a unique fixed point
there.  The inequality $e<1$ also implies that $AJ_F(c)$ is invertible,
hence $A$ is invertible.  Therefore the fixed point satisfies $F(x)=0$,
not merely $AF(x)=0$.  This proves exact existence of the claimed
orthonormal basis and all its AC identities.

\section{Verification and provenance}
\texttt{verify.py} uses only standard-library integer and rational
arithmetic.  It reconstructs the complete moment system, recomputes the
Jacobian, the inverse defect, the correction bound, and the quadratic
Lipschitz bound, then checks every interval and coefficient association
in \texttt{basis.json}.  The separate independent verifier reconstructs
the same conditions from physical spin matrices.

The numerical seed is the supplied Julia subspace
\(\Pi\)\texttt{\_20\_4\_2}.  A numerical refinement suggests the
certificate center; the rational contraction proof establishes the
result.  Equality to the original decimal span is not asserted.
No claim of literature priority is made.
\end{document}
